\section{Theory}

The measurement problem arises because an LGFV is both a legal organization and
a carrier of local fiscal functions. A change in the organization can leave the
function intact, while a change in the function can occur without liquidation
of the organization. The theory therefore begins with the location of the
function and then asks which local conditions are associated with different
institutional fates.

\subsection{Official and institutional exit}

Official exit is a documented administrative or corporate event. A platform is
removed from a list, declares that it no longer performs government financing,
is described as market-oriented, or undergoes restructuring or termination.
Institutional exit concerns the financing, project, and liability arrangements
that follow. The concepts are related because the formal event can initiate a
real reorganization, but they are not equivalent. Regulatory language identifies
the event to be investigated; it does not determine the outcome.

This distinction follows from the fiscal role of LGFVs. They link local
governments to infrastructure, land development, credit, and state assets
\citep{wong2013,tsui2011,jin_rial_2016_lgfv_ppp}. A firm can cease issuing debt
on behalf of the government while continuing fee-based project management
funded through the budget. It can also declare commercial transformation while
retaining entrusted construction, fiscal receivables, land preparation, or
government-linked repayment. A third possibility is that the original firm
changes role while another local public entity assumes the same function. The
post-event location of the function separates these outcomes.

\subsection{Institutional fates}

Substantive exit occurs when the issuer stops performing the relevant
government financing or quasi-fiscal function. The company may continue as a
commercial firm, and public projects may continue through budgetary spending,
special-purpose bonds, or fee-based management, but the platform no longer
serves as an off-budget fiscal arm. Nominal exit occurs when formal compliance
language appears while the same issuer continues the relevant function.
Functional transfer occurs when the original issuer exits or is reorganized and
another local state-owned entity, subsidiary, or merged group assumes the
function. Liquidation occurs when the legal entity terminates. It is coded
separately because termination does not establish whether the fiscal function
ended or migrated.

These categories distinguish policy-practice decoupling from organizational
substitution. Nominal exit is a gap between formal status and the same issuer's
practice. Functional transfer can involve genuine change in that issuer while
leaving the broader local fiscal arrangement intact. Substantive exit requires
evidence that the financing function itself no longer resides in a platform-like
entity. Liquidation records an organizational endpoint and therefore requires a
separate search for any successor function.

\subsection{Contemporary processes}

Three contemporary processes can produce movement away from nominal exit.
Fiscal absorption occurs when an eligible liability or project cost enters an
approved budget, appropriation, statutory debt register, or government-bond
allocation and quasi-fiscal rollover falls. Debt conversion alone is
insufficient because financing may migrate to another issuer or off-budget
instrument \citep{wingender_2018_intergovernmental,
chen_he_liu_2020_financing}. Positive evidence must identify the public payer
and the legal financing instrument after the event.

Bureaucratic coordination occurs when finance bureaus, development and reform
agencies, state-owned asset supervisors, creditors, and issuers make consistent
assignments and complete the transfer of liabilities, projects, staff, data, or
monitoring responsibility. Uniform documents do not establish this process.
Coordination can sustain evasion as well as implementation, so the evidence must
show completed assignments and follow-up rather than a common statement alone
\citep{zhou_2010_collusion}.

Project or state-asset governance occurs when post-event records clarify title,
valuation, contracts, liabilities, operating rights, personnel, and revenue
control for transferred projects or assets. Regrouping or a debt swap without
changes in control does not establish it. This process differs from fiscal
capacity because a locality can have ample resources while retaining weak
boundaries between public projects, state-owned enterprise activity, and
implicit government financing.

The three processes are analytically distinct. A subsidy can identify a payer
without showing cross-agency implementation. A transfer decision can show
coordination without establishing who controls the asset after the transfer. A
new project title can show legal control without showing that a liability moved
onto the budget. Each process therefore requires its own positive evidence,
falsifier, source-quality field, and missing-evidence state.

\subsection{Historical capacity as an exploratory explanation}

Historical administrative density and elite formation may be associated with
contemporary institutions, but the available historical measure is not a direct
observation of fiscal absorption, coordination, or asset governance. The
empirical application uses Ming-Qing examination-elite density matched to
contemporary prefectures. This measure can capture long-run administrative
human capital and state-society ties, but it also predicts modern education,
culture, and social capital and does not establish continuity in political
elites \citep{chen_kung_ma_2020_keju}. Calling it historical state capacity is
therefore an interpretation that must be tested against those alternatives.

The application asks whether places with higher historical elite density are
more likely to move away from nominal exit. One possible interpretation is that
long-run administrative development left routines and organizations that make
formal fiscal substitution easier. Other interpretations include contemporary
wealth, educational infrastructure, social capital, provincial policy,
document availability, and the composition of platform issuers. The present
design does not distinguish these channels or identify historical mediation.

\subsection{Empirical expectations}

The first expectation is associational: higher historical elite density should
coincide with a greater probability of movement away from nominal exit. The
second is institutional: cases coded as substantive exit should contain more
evidence of fiscal absorption, while functional-transfer cases should contain
more evidence of completed reassignment and post-transfer governance. Nominal
exit should contain continued-function evidence for the same issuer and fail
the principal falsifiers of those processes. These expectations discipline the
case evidence, but only independently coded process indicators and a stronger
comparative design could support a mechanism or causal-persistence claim.
